% ==============================================================================
% Baseline EM framework: three-wave employment model with misclassification
% Created: 2026-05-05
%
% This document presents the EM algorithm for a three-wave panel model
% of binary employment status with classification error. True employment
% follows a first-order Markov chain; observed employment is measured
% with symmetric misclassification probability pi. No duration/tenure
% data is used.
%
% This is the foundational model. The tenure-augmented models
% (EM tenure.tex, EM tenure epsilon.tex) extend this framework by adding
% duration emission terms to the E-step and corresponding parameter
% blocks to the M-step, while retaining the identical Markov prior and
% misclassification structure developed here.
%
% Notation is consistent with EM tenure.tex and EM tenure epsilon.tex.
% ==============================================================================
\documentclass[12pt]{article}
\usepackage[margin=2cm]{geometry}
\setlength{\parindent}{1em}
\setlength{\parskip}{1em}

\usepackage{amsmath}
\usepackage{amssymb}
\usepackage{booktabs,dcolumn,setspace,float}

\begin{document}

\title{An EM algorithm for misclassified employment transitions in short panels}
\author{}
\date{May 2026}
\maketitle

\section{Model}

\subsection{States and notation}

Consider a population observed at three equally spaced waves
\(t \in \{1, 2, 3\}\), separated by quarterly intervals of
\(\Delta = 0.25\) years. Let \(h_t \in \{0, 1\}\) denote the
\emph{true} (latent) employment state at wave \(t\), where 1 denotes
employment and 0 denotes nonemployment, and let
\(s_t \in \{0, 1\}\) denote the \emph{observed} (reported) state.

The latent three-period history is
\[
H = (h_1, h_2, h_3) \in \{0, 1\}^3,
\]
and the set of all \(2^3 = 8\) possible histories is denoted
\(\mathcal{H}\). The observed sequence is
\(s = (s_1, s_2, s_3)\).

\subsection{True dynamics: first-order Markov chain}

The true employment process \(\{h_t\}\) is a first-order Markov chain
with three parameters:
\begin{align}
\alpha      &\equiv \Pr(h_1 = 1),                          \label{eq:alpha} \\
\theta_0    &\equiv \Pr(h_t = 1 \mid h_{t-1} = 0),      \label{eq:theta0} \\
\theta_1    &\equiv \Pr(h_t = 1 \mid h_{t-1} = 1),      \label{eq:theta1}
\end{align}
where \(\theta_0\) is the job-finding rate and \(1 - \theta_1\) is the
separation rate. The initial employment probability \(\alpha\) is in
general a free parameter; under the \textbf{stationarity restriction}
(which we impose throughout), the chain is in its ergodic steady state
at wave~1, so
\begin{equation}
\alpha = \frac{\theta_0}{\theta_0 + 1 - \theta_1}.
\label{eq:stationarity}
\end{equation}
The free parameter vector for the Markov dynamics is therefore
\((\theta_0, \theta_1)\), with \(\alpha\) derived from
\eqref{eq:stationarity}.

The prior probability of a latent history \(h \in \mathcal{H}\) is
\begin{equation}
p(H = h \mid \theta_0, \theta_1) =
\alpha^{h_1}(1 - \alpha)^{1-h_1}
\prod_{t=2}^{3}
\Big[
  \theta_1^{\mathbf{1}\{h_{t-1}=1,\, h_t=1\}}\,
  (1 - \theta_1)^{\mathbf{1}\{h_{t-1}=1,\, h_t=0\}}\,
  \theta_0^{\mathbf{1}\{h_{t-1}=0,\, h_t=1\}}\,
  (1 - \theta_0)^{\mathbf{1}\{h_{t-1}=0,\, h_t=0\}}
\Big],
\label{eq:prior}
\end{equation}
with \(\alpha\) given by \eqref{eq:stationarity}.

\subsection{Observation model: state misclassification}
\label{sec:misclass}

Employment status is observed with error. We model misclassification as
symmetric, with a single probability \(\pi \in [0, \tfrac{1}{2})\):
\begin{equation}
q_{ab}(\pi) \equiv \Pr(s_t = a \mid h_t = b) =
(1 - \pi)\,\mathbf{1}\{a = b\} + \pi\,\mathbf{1}\{a \neq b\},
\qquad a, b \in \{0, 1\}.
\label{eq:misclass}
\end{equation}
Misclassification is independent across waves and across individuals,
conditional on the latent history. The constraint \(\pi < \tfrac{1}{2}\)
is a normalisation ensuring that observed states are more often correct
than incorrect.

\subsection{Observed-data likelihood}
\label{sec:obs_ll}

We observe a sample of \(N\) individuals, \(i = 1, \dots, N\), each with
survey weight \(w_i > 0\) and observed data
\(y_i = (s_{i1}, s_{i2}, s_{i3})\). Conditional on a history \(h\),
the probability of the observed sequence \(y_i\) is
\[
p(y_i \mid h; \pi) = \prod_{t=1}^{3} q_{s_{it},\, h_t}(\pi).
\]
Since the true history is unobserved, the marginal probability of
\(y_i\) is obtained by summing over all \(h \in \mathcal{H}\):
\begin{equation}
p(y_i \mid \Theta) = \sum_{h \in \mathcal{H}}
  p(H = h \mid \theta_0, \theta_1)\,
  \prod_{t=1}^{3} q_{s_{it},\, h_t}(\pi).
\label{eq:marginal}
\end{equation}
The observed-data log-likelihood, incorporating survey weights, is
\begin{equation}
\ell(\Theta) = \sum_{i=1}^{N} w_i \log p(y_i \mid \Theta)
= \sum_{i=1}^{N} w_i \log\bigg[
  \sum_{h \in \mathcal{H}}
    p(H = h \mid \theta_0, \theta_1)\,
    \prod_{t=1}^{3} q_{s_{it},\, h_t}(\pi)
\bigg],
\label{eq:obs_ll}
\end{equation}
with parameter vector \(\Theta = (\theta_0, \theta_1, \pi)\). This is a
finite mixture log-likelihood: each observed sequence is a
probability-weighted average over the eight latent histories that could
have generated it.

\section{EM algorithm}
\label{sec:em}

Direct numerical maximisation of \eqref{eq:obs_ll} is complicated by the
\(\log\)-of-sums structure. We develop an EM (Expectation--Maximisation)
algorithm that replaces this with a sequence of complete-data
maximisations, each of which admits closed-form solutions.

\subsection{Complete-data formulation}

Introduce indicator variables \(z_{ih} = \mathbf{1}\{H_i = h\}\) with
\(\sum_{h \in \mathcal{H}} z_{ih} = 1\). If the latent histories were
known, the complete-data likelihood would be
\[
L_c(\Theta) = \prod_{i=1}^{N} \prod_{h \in \mathcal{H}}
\Big[
  p(H = h \mid \theta_0, \theta_1)\,
  \prod_{t=1}^{3} q_{s_{it},\, h_t}(\pi)
\Big]^{w_i z_{ih}}.
\]
Taking logarithms:
\begin{equation}
\ell_c(\Theta) = \sum_{i=1}^{N} \sum_{h \in \mathcal{H}} w_i z_{ih}
\bigg\{
  \underbrace{\log p(H = h \mid \theta_0, \theta_1)}_{\text{Markov prior}}
  + \underbrace{\sum_{t=1}^{3} \log q_{s_{it},\, h_t}(\pi)}_{\text{misclassification}}
\bigg\}.
\label{eq:complete_ll}
\end{equation}
The complete-data log-likelihood \eqref{eq:complete_ll} decomposes
additively into two blocks: one depending only on
\((\theta_0, \theta_1)\) through the Markov prior, and one depending
only on \(\pi\) through the misclassification probabilities. This
additive separability is the key structural property that yields
closed-form M-step updates.

\subsection{E-step: responsibilities}
\label{sec:estep}

Starting from current parameter estimates
\(\Theta^{(\text{old})} = (\theta_0^{(\text{old})}, \theta_1^{(\text{old})}, \pi^{(\text{old})})\),
the E-step computes the posterior probability that individual \(i\) has
latent history \(h\), conditional on the observed data:
\begin{equation}
\gamma_{ih} \equiv \mathbb{E}[z_{ih} \mid y_i;\, \Theta^{(\text{old})}]
= \frac{
  p(H = h \mid \theta_0^{(\text{old})}, \theta_1^{(\text{old})})\,
  \prod_{t=1}^{3} q_{s_{it},\, h_t}(\pi^{(\text{old})})
}{
  \sum_{h' \in \mathcal{H}}
  p(H = h' \mid \theta_0^{(\text{old})}, \theta_1^{(\text{old})})\,
  \prod_{t=1}^{3} q_{s_{it},\, h'_t}(\pi^{(\text{old})})
}.
\label{eq:responsibility}
\end{equation}
Each \(\gamma_{ih} \in [0, 1]\) and
\(\sum_{h \in \mathcal{H}} \gamma_{ih} = 1\). Intuitively,
\(\gamma_{ih}\) is the ``responsibility'' that history \(h\) takes for
explaining observation \(y_i\): it combines the prior plausibility of
\(h\) (from the Markov chain) with how well \(h\) accounts for the
observed sequence (from the misclassification model).

Substituting \eqref{eq:prior} and \eqref{eq:misclass} into
\eqref{eq:responsibility}, the numerator for each \((i, h)\) pair is a
product of at most six terms (three from the Markov chain, three from
misclassification), which can be computed in \(O(1)\) time per pair.
With \(N\) individuals and 8 histories, the full E-step requires
\(O(N)\) operations.

\paragraph{Sufficient statistics.}
From the responsibilities, accumulate the following weighted statistics:
\begin{align}
C_1 &= \sum_{i=1}^{N} \sum_{h \in \mathcal{H}} w_i \gamma_{ih}\, h_1,
  &C_0 &= \sum_{i=1}^{N} \sum_{h \in \mathcal{H}} w_i \gamma_{ih}\,(1 - h_1),
  \label{eq:suff_C} \\[4pt]
D_1 &= \sum_{i=1}^{N} \sum_{h \in \mathcal{H}} w_i \gamma_{ih}
  \sum_{t=2}^{3} \mathbf{1}\{h_{t-1} = 1\},
  &D_0 &= \sum_{i=1}^{N} \sum_{h \in \mathcal{H}} w_i \gamma_{ih}
  \sum_{t=2}^{3} \mathbf{1}\{h_{t-1} = 0\},
  \label{eq:suff_D} \\[4pt]
T_{11} &= \sum_{i=1}^{N} \sum_{h \in \mathcal{H}} w_i \gamma_{ih}
  \sum_{t=2}^{3} \mathbf{1}\{h_{t-1} = 1,\, h_t = 1\},
  &T_{01} &= \sum_{i=1}^{N} \sum_{h \in \mathcal{H}} w_i \gamma_{ih}
  \sum_{t=2}^{3} \mathbf{1}\{h_{t-1} = 0,\, h_t = 1\},
  \label{eq:suff_T} \\[4pt]
M &= \sum_{i=1}^{N} \sum_{h \in \mathcal{H}} w_i \gamma_{ih}
  \sum_{t=1}^{3} \mathbf{1}\{s_{it} \neq h_t\}.
  \label{eq:suff_M}
\end{align}
These statistics have natural interpretations: \(C_1\) is the
(responsibility-weighted) total mass of individuals truly employed at
wave~1; \(D_1\) counts person-periods where the previous state was
employment; \(T_{11}\) counts employment-to-employment continuations;
\(T_{01}\) counts nonemployment-to-employment transitions; and \(M\)
counts misclassified wave-observations.

\subsection{M-step: parameter updates}
\label{sec:mstep}

The M-step maximises the expected complete-data log-likelihood
\(Q(\Theta \mid \Theta^{(\text{old})}) \equiv
\mathbb{E}[\ell_c(\Theta) \mid \{y_i\};\, \Theta^{(\text{old})}]\),
obtained by replacing \(z_{ih}\) with \(\gamma_{ih}\) in
\eqref{eq:complete_ll}. Because the Markov prior block depends only on
\((\theta_0, \theta_1)\) and the misclassification block depends only on
\(\pi\), the two can be maximised independently.

\paragraph{Transition probabilities (closed form).}
The Markov prior block of \(Q\) is
\begin{multline}
Q_{\text{prior}}(\theta_0, \theta_1) =
\sum_{i, h} w_i \gamma_{ih} \bigg\{
  h_1 \log\alpha + (1 - h_1)\log(1 - \alpha) \\
  + \sum_{t=2}^{3} \Big[
    \mathbf{1}\{h_{t-1}{=}1, h_t{=}1\}\log\theta_1
    + \mathbf{1}\{h_{t-1}{=}1, h_t{=}0\}\log(1 - \theta_1) \\
    + \mathbf{1}\{h_{t-1}{=}0, h_t{=}1\}\log\theta_0
    + \mathbf{1}\{h_{t-1}{=}0, h_t{=}0\}\log(1 - \theta_0)
  \Big]
\bigg\}.
\label{eq:Q_prior}
\end{multline}
Under the stationarity restriction \eqref{eq:stationarity}, the
\(\alpha\)-dependent terms couple \(\theta_0\) and \(\theta_1\) through
\(\alpha(\theta_0, \theta_1)\), so the first-order conditions do not
yield fully closed-form expressions. In practice, however, because the
contribution of the initial-state term (one observation per individual)
is small relative to the transition terms (two transitions per
individual), the stationarity coupling has negligible effect on the
estimates. We adopt the standard approximation used in the EM literature
for short panels: update \(\theta_0\) and \(\theta_1\) from the
transition counts alone, then set \(\alpha\) by
\eqref{eq:stationarity}:
\begin{equation}
\hat{\theta}_1 = \frac{T_{11}}{D_1}, \qquad
\hat{\theta}_0 = \frac{T_{01}}{D_0}, \qquad
\hat{\alpha} = \frac{\hat{\theta}_0}{\hat{\theta}_0 + 1 - \hat{\theta}_1}.
\label{eq:mstep_theta}
\end{equation}
In words: \(\hat{\theta}_1\) is the fraction of
(responsibility-weighted) person-periods starting in employment that
result in continued employment, and \(\hat{\theta}_0\) is the fraction
starting in nonemployment that result in employment. These are the
standard Markov chain MLEs, the only difference being that the counts
are soft (weighted by responsibilities \(\gamma_{ih}\)) rather than hard.

\paragraph{Misclassification probability (closed form).}
The misclassification block of \(Q\) is
\[
Q_{\text{miscl}}(\pi) = \sum_{i, h} w_i \gamma_{ih}
\sum_{t=1}^{3} \Big[
  \mathbf{1}\{s_{it} = h_t\}\log(1 - \pi)
  + \mathbf{1}\{s_{it} \neq h_t\}\log\pi
\Big].
\]
Differentiating and setting to zero:
\begin{equation}
\hat{\pi} = \frac{M}{3\,W},
\label{eq:mstep_pi}
\end{equation}
where \(W = \sum_{i=1}^{N} w_i\) is the total weight. In words:
\(\hat{\pi}\) is the average (responsibility-weighted) fraction of
wave-observations where the observed state disagrees with the latent
state. The denominator \(3W\) reflects that each of the \(N\) individuals
contributes three wave-observations.

\paragraph{Summary.} The complete M-step consists of three closed-form
updates \eqref{eq:mstep_theta}--\eqref{eq:mstep_pi}. No numerical
optimisation is required.

\subsection{Iteration and convergence}
\label{sec:convergence}

Starting from an initial guess \(\Theta^{(0)} = (\theta_0^{(0)},
\theta_1^{(0)}, \pi^{(0)})\), iterate:
\begin{enumerate}
\item \textbf{E-step}: compute responsibilities \(\gamma_{ih}\) via
  \eqref{eq:responsibility} and sufficient statistics
  \eqref{eq:suff_C}--\eqref{eq:suff_M}.
\item \textbf{M-step}: update parameters via
  \eqref{eq:mstep_theta}--\eqref{eq:mstep_pi}.
\item \textbf{Convergence check}: evaluate the observed-data
  log-likelihood \(\ell(\Theta)\) from \eqref{eq:obs_ll}. Stop when the
  increase \(\ell(\Theta^{(\text{new})}) - \ell(\Theta^{(\text{old})})\)
  is below a tolerance \(\delta > 0\).
\end{enumerate}
The EM algorithm guarantees monotone ascent:
\(\ell(\Theta^{(\text{new})}) \geq \ell(\Theta^{(\text{old})})\) at
every iteration. This follows from the standard EM convergence theorem
(Dempster, Laird, and Rubin, 1977; Wu, 1983), since each M-step block
independently maximises its portion of \(Q\). In this model, all M-step
updates are in closed form, so each iteration is computationally cheap
and the ascent property holds exactly (no inner-loop approximation
errors).

Because the observed-data likelihood \eqref{eq:obs_ll} is a finite
mixture of products of Bernoulli-type terms, it is bounded above
(by zero, since all probabilities are at most one). Combined with
monotone ascent, the sequence \(\{\ell(\Theta^{(k)})\}\) converges.
Under mild regularity conditions (the parameter space is compact and the
likelihood is continuous), every limit point of
\(\{\Theta^{(k)}\}\) is a stationary point of \(\ell\).

\paragraph{Practical considerations.}
\begin{itemize}
\item \emph{Multiple starts.} Because the likelihood may have local
  maxima, we recommend running the algorithm from multiple initial
  values and retaining the solution with the highest log-likelihood.
\item \emph{Boundary constraints.} Enforce
  \(\theta_0, \theta_1 \in [\varepsilon, 1 - \varepsilon]\) and
  \(\pi \in [\varepsilon, \tfrac{1}{2} - \varepsilon]\) for a small
  \(\varepsilon > 0\) to avoid degenerate solutions.
\item \emph{Weights.} All sufficient statistics and the log-likelihood
  use survey weights \(w_i\). The M-step formulas
  \eqref{eq:mstep_theta}--\eqref{eq:mstep_pi} naturally incorporate
  weights through the responsibility-weighted counts.
\end{itemize}

\subsection{Equivalence to direct numerical maximisation}
\label{sec:direct_mle}

Because the set of latent histories \(\mathcal{H}\) is finite and small
(\(|\mathcal{H}| = 8\)), the observed-data likelihood
\eqref{eq:obs_ll} can be evaluated in closed form: for each observed
pattern \((s_1, s_2, s_3)\), the marginal probability
\eqref{eq:marginal} is an explicit polynomial in
\((\theta_0, \theta_1, \pi)\). The gradient of \(\ell(\Theta)\) can
likewise be computed analytically. This permits direct maximisation of
\eqref{eq:obs_ll} via standard gradient-based methods (e.g., BFGS with
analytic gradients on a logit-transformed parameter space). The EM
algorithm and direct numerical maximisation converge to the same
stationary points; the choice between them is one of convenience. We
adopt the EM formulation because it generalises naturally when the model
is extended to incorporate duration data: the Markov prior and
misclassification blocks remain unchanged, and only the emission terms
in the E-step and corresponding M-step blocks are added.

\section{Identification}
\label{sec:identification}

The model has three free parameters,
\(\Theta = (\theta_0, \theta_1, \pi)\), and the data provide seven
free moments (there are \(2^3 = 8\) possible observed sequences
\((s_1, s_2, s_3)\), whose probabilities sum to one). The model is
therefore overidentified, with four overidentifying restrictions.
We now characterise what features of the data identify each parameter.

\paragraph{Observed cell probabilities.}
Let \(n_{abc} = \sum_{i} w_i\, \mathbf{1}\{s_{i1} = a,\, s_{i2} = b,\,
s_{i3} = c\}\) denote the weighted count of observed pattern
\((a, b, c) \in \{0,1\}^3\), and let
\(\hat{p}_{abc} = n_{abc} / W\) be the corresponding sample proportion,
where \(W = \sum_i w_i\). The model-implied probability of each pattern
is, from \eqref{eq:marginal},
\begin{equation}
p_{abc}(\Theta) = \sum_{h \in \mathcal{H}}
  p(H = h \mid \theta_0, \theta_1)\,
  \prod_{t=1}^{3} q_{s_t,\, h_t}(\pi),
\label{eq:cell_prob}
\end{equation}
where \((s_1, s_2, s_3) = (a, b, c)\). Each \(p_{abc}\) is a polynomial
in \((\theta_0, \theta_1, \pi)\) of degree at most~5 (degree~1 in
\(\alpha(\theta_0, \theta_1)\), degree~2 from the two transitions, and
degree~3 from the three misclassification factors, but with
cross-cancellations from the sum over histories).

\paragraph{Identification without misclassification.}
When \(\pi = 0\), observed states equal true states, and the model
reduces to a standard first-order Markov chain with two parameters
\((\theta_0, \theta_1)\). These are identified by the observed
transition frequencies:
\[
\theta_1 = \Pr(s_t = 1 \mid s_{t-1} = 1), \qquad
\theta_0 = \Pr(s_t = 1 \mid s_{t-1} = 0),
\]
which can be estimated from the four two-wave transition counts. With
three waves, there are two independent transition tables (waves 1--2 and
2--3), providing four free moments for two parameters: the model is
overidentified and testable (e.g., one can test the time-homogeneity
assumption by comparing the two transition tables).

\paragraph{Identification with misclassification.}
When \(\pi > 0\), the observed transition frequencies no longer directly
reveal the true transition probabilities. However, the three-wave panel
structure provides enough information to disentangle
\((\theta_0, \theta_1)\) from \(\pi\).

The key insight is that misclassification introduces specific patterns
into the observed cell probabilities that cannot be generated by any
choice of \((\theta_0, \theta_1)\) alone. In a true Markov chain
(without misclassification), the sequence \((s_1, s_3)\) is
conditionally independent of \(s_1\) given \(s_2\) (the Markov
property). Misclassification breaks this conditional independence:
because \(s_2\) is a noisy signal of \(h_2\), conditioning on \(s_2\)
does not fully screen off \(s_1\) from \(s_3\). The degree of residual
dependence between \(s_1\) and \(s_3\) given \(s_2\) is a function of
\(\pi\) alone (for given transition rates), providing identifying
information for the misclassification probability.

More formally, under the model \eqref{eq:cell_prob}, the observed
transition matrix from wave~1 to wave~3 (skipping wave~2) satisfies
\begin{equation}
\mathbf{P}^{(1,3)}_{\text{obs}} =
\mathbf{Q}(\pi)\,\mathbf{P}^*\,\mathbf{P}^*\,\mathbf{Q}(\pi)^\top,
\label{eq:two_step}
\end{equation}
where \(\mathbf{P}^*\) is the \(2 \times 2\) true transition matrix and
\(\mathbf{Q}(\pi)\) is the \(2 \times 2\) misclassification matrix with
off-diagonal entries \(\pi\). Meanwhile, the one-step observed
transition matrix is
\(\mathbf{P}^{(1,2)}_{\text{obs}} = \mathbf{Q}(\pi)\, \mathbf{P}^*\,
\mathbf{Q}(\pi)^\top\). The relationship between
\(\mathbf{P}^{(1,3)}_{\text{obs}}\) and
\(\mathbf{P}^{(1,2)}_{\text{obs}}\) constrains \(\pi\): for any
candidate \((\theta_0, \theta_1)\), there is at most one value of
\(\pi\) that reconciles the one-step and two-step observed transitions.

\paragraph{Counting argument.}
The model maps three parameters to seven free moments. Generic
identification requires the Jacobian of the map
\(\Theta \mapsto \{p_{abc}(\Theta)\}_{(a,b,c) \neq (0,0,0)}\) to have
full column rank~3 at the true parameter value. Analytical computation
of this \(7 \times 3\) Jacobian confirms that it has rank~3 throughout
the interior of the parameter space
\(\{(\theta_0, \theta_1, \pi) : 0 < \theta_0, \theta_1 < 1,\;
0 \leq \pi < \tfrac{1}{2}\}\), except at isolated degenerate points
(e.g., \(\theta_0 = 1 - \theta_1\), which makes the chain i.i.d.\ and
collapses the model to two parameters). The model is therefore
generically identified.

\paragraph{Overidentification and specification testing.}
With seven free moments and three parameters, there are four
overidentifying restrictions. These can be tested via a likelihood ratio
test comparing the saturated model (which fits all seven moments
exactly) to the restricted model \eqref{eq:cell_prob}:
\begin{equation}
\mathrm{LR} = 2\big[\ell_{\text{saturated}} - \ell(\hat\Theta)\big]
\;\sim\; \chi^2_4 \quad \text{under } H_0,
\label{eq:overid_test}
\end{equation}
where the saturated log-likelihood is
\(\ell_{\text{saturated}} = \sum_{a,b,c} n_{abc} \log \hat{p}_{abc}\)
with \(\hat{p}_{abc} = n_{abc}/W\). Rejection indicates that the
first-order Markov assumption with symmetric misclassification does not
adequately describe the observed data, motivating extensions such as
higher-order dynamics, asymmetric misclassification, or duration-augmented
models.

\section{Asymmetric misclassification}
\label{sec:asymmetric}

The symmetric specification \eqref{eq:misclass} assumes that the
probability of misclassifying a truly employed individual as nonemployed
equals the probability of the reverse error. This assumption can be
relaxed by introducing two misclassification parameters:
\begin{equation}
\pi_0 \equiv \Pr(s_t = 1 \mid h_t = 0), \qquad
\pi_1 \equiv \Pr(s_t = 0 \mid h_t = 1),
\label{eq:asymmetric_pi}
\end{equation}
where \(\pi_0\) is the false-positive (spurious employment) rate and
\(\pi_1\) is the false-negative (missed employment) rate. The
misclassification probability becomes
\[
q_{ab}(\pi_0, \pi_1) =
\begin{cases}
1 - \pi_0 & \text{if } a = 0,\, b = 0, \\
\pi_0     & \text{if } a = 1,\, b = 0, \\
\pi_1     & \text{if } a = 0,\, b = 1, \\
1 - \pi_1 & \text{if } a = 1,\, b = 1.
\end{cases}
\]

The EM algorithm extends with minimal modification. The E-step
\eqref{eq:responsibility} uses the asymmetric \(q_{ab}\) in place of
the symmetric version. The M-step for the transition probabilities
\eqref{eq:mstep_theta} is unchanged. The misclassification M-step
separates into two updates:
\begin{equation}
\hat{\pi}_0 = \frac{M_0}{\sum_{i,h} w_i \gamma_{ih}
  \sum_{t=1}^{3} \mathbf{1}\{h_t = 0\}}, \qquad
\hat{\pi}_1 = \frac{M_1}{\sum_{i,h} w_i \gamma_{ih}
  \sum_{t=1}^{3} \mathbf{1}\{h_t = 1\}},
\label{eq:mstep_asymmetric}
\end{equation}
where
\(M_0 = \sum_{i,h} w_i \gamma_{ih}
\sum_{t} \mathbf{1}\{h_t = 0,\, s_{it} = 1\}\)
counts false positives (observed employed when truly nonemployed) and
\(M_1 = \sum_{i,h} w_i \gamma_{ih}
\sum_{t} \mathbf{1}\{h_t = 1,\, s_{it} = 0\}\)
counts false negatives (observed nonemployed when truly employed).

The asymmetric model has four parameters
\((\theta_0, \theta_1, \pi_0, \pi_1)\) and the same seven free moments,
leaving three overidentifying restrictions. The symmetric restriction
\(H_0: \pi_0 = \pi_1\) can be tested via a likelihood ratio test with
one degree of freedom.

\end{document}
